\section{The three indices preceding the uniform local range}
\label{ext:floor40}

The local argument at $n=40,41,42$ uses the same contact, Fourier and
conditional-count identities as Section~\ref{ext:local-engine}. There are
two extra issues. A bound computed with index $40$ must remain valid at
the actual index, and a signed correction divided by $n$ cannot always
be enlarged by replacing $n$ by $40$.

Suppose that $40\le n\le42$ is a first strict failure above a sharp base
through $39$. The recurrence of Theorem~\ref{ext:one-step} gives
\[
 0<n(C_n-\CE)<\frac{40}{39}
 \left(\frac{41}{200}+\frac{17}{250}+\frac{17}{1250}\right)
 =\frac{1433}{4875}<\frac3{10},\qquad C_n<\frac{167}{400}<.4226.
\]
Here $\sqrt{n/(n-1)}<51/50$ and $\sqrt{n-1}>6$ were used.
Both deleted constants are at most $\CE<.4253$ and the predecessor loss
is zero. In particular all bounds remain valid if the local estimates
use the weaker excess allowance $.54/n$ and deletion allowance $.19$.

\subsection{Why a computation at floor 40 can be reused}

The invariant on one fixed actual law is
\begin{equation}\label{ext:floor40-invariant}
 \mathcal R_{\mathrm{actual}}\le r,\qquad nD_{\mathrm{actual}}\le40d.
\end{equation}
The second inequality implies $D_{\mathrm{actual}}\le d$, which may be
used in index-free moment and endpoint bounds. Its converse is not used.
Every contact update bounds $\kappa nD$ by a phase budget $B$, where
$\kappa=997/15000$, and therefore exports $d=B/(40\kappa)$.
A negative budget excludes the law. Taking minima with established
bounds or maxima over a complete cover preserves this invariant.

For example, suppose the joint resource gives
\[
 D\ge b+6(1-\theta)(A+B)d_{\rm in}+K\theta,
 \qquad b,K\ge0,
\]
and $\theta\le\Theta<1$, $pq\le v_+$. Since
$\lambda=npq\theta/(1-\theta)$, multiplication by $n$ yields
\[
 40(d-b)\ge40(1-\Theta)6(A+B)d_{\rm in}
            +\frac{K(1-\Theta)}{v_+}\lambda.
\]
This proves the saved joint intercepts for the actual noise $\lambda$.
On a noise band $[a,b']$, the further simultaneous restrictions are
\[
 d_{\rm in}\le d_0(1-a/\Lambda),\qquad
 \theta\le\min\left\{\Theta,\frac{b'}{40v_-+b'}\right\}.
\]
No different law is selected when one of these bounds improves another.
If a saved harmonic geometry bound has the form
\[
 \theta\le\left(1+\frac{nv_-}{\Lambda}
                   +\frac{v_-}{B_0d_0}\right)^{-1},
 \qquad B_0,d_0,\Lambda>0,
\]
its value at $40$ is an upper bound. This substitution bounds $\theta$;
the actual count noise remains $\lambda=npq\theta/(1-\theta)$.

The nonnegative Fourier error coefficients decrease with the actual
index. The low-frequency factors are products and sums of
$(n/(n-1))^{3/2}$, $(n-1)^{-1}$, $n^{-1/2}$ and inverse powers of
$(n-1)/(2n)$ and $(45-M)(1-2/n)/96$. The latter bases are positive
and increasing on their retained domains. For the high-frequency
terms, $\sqrt n e^{-cn}$ decreases when $2\cdot40c>1$; the other
factors are $n^{-1}$, $e^{-knb^2}/\sqrt n$ and the positive decreasing
Mills bound $(x+2)/(x^2+4x+2)$.
The optional product $\sqrt x e^{-x}(x+2)/(x^2+4x+2)$ decreases for
$x\ge1/4$: the numerator of its logarithmic derivative, after
multiplication by a positive denominator, is
$4-10x-22x^2-13x^3-2x^4<0$ there. Each use retains its executable
positivity and domain gates.

The density lower bound $\phi(0)(1-5r/80)$ remains valid later.
Signed coefficients are minimized over the whole density interval;
in particular the paired curvature is
\[
 C\rho+\frac12\min\{\phi_lg_0,(2/5)g_0\},
 \qquad g_0=(79/25)\rho-3-d.
\]
This is increasing with $\phi_l$ if $g_0\ge0$, and constant at the
upper-density endpoint otherwise. This verifies the sign-dependent
direction of the floor substitution. The positive conserved phase
price and both signed linear endpoints are retained.
The general, unpaired curvature has lower-density derivative
$(g_0+(d-4/25)_+)/2\ge0$: for $d\le4/25$ one has
$g_0\ge4/25-d$, and otherwise the numerator is
$(79/25)(\rho-1)$. The conserved addition has coefficient
$(.4095\rho-d/15)(1-5r/80)/2>0$ on each admitted cell.
The own-score numerator
\[
 \frac r{40}-\phi_{0,l}(1-5r/80)\frac{D_l}{6}
\]
also bounds its actual-index counterpart when $r,D_l\ge0$.
For a strict failure, the additional sharp-score charge satisfies
\[
 n(C-\CE)\rho\ge40\max\{0,C-.40974\}\rho_l.
\]
Indeed $\CE<.40974$, and each replacement decreases a nonnegative
factor on the part where $C>.40974$. The score cover includes that
breakpoint. The charge is affine on either side; hence the coherent
$(C,\rho)$ corner argument still applies before maximizing over the
threshold. These statements concern the zero-barrier sharp target.

The first support proof is evaluated separately on the full band
$40\le n\le42$, with zero loss and radius $96/25$. Its favourable
fifth-derivative term uses the upper endpoint $42$; the value at index
$40$ alone would not suffice. Every later left radius contraction
also carries this two-endpoint proof. In the central contact expansion,
all index-dependent error numerators are nonnegative after taking
positive parts, so their values at $40$ bound their later values.
The already acquired support closes the exterior of every such expansion.
In particular its complete cell bound is
\[
 L_P+\frac{.54Q_+}{n}
 +.4\left(\sum_{j=1}^5\frac{M_j}{n^j}+\frac{T}{n^6}\right)
 +\rho\left(\sqrt{\frac n{n-1}}\delta_n+
                          \frac{.384(.4226)}n\right),
\]
where $Q_+,M_j,T\ge0$, $L_P$ is independent of $n$ and may be
negative, and $\delta_n\le.19$. Every positive term is bounded by
its value at $40$. This discards favourable Gaussian terms and is
therefore distinct from the two-endpoint initial support polynomial.

For concentration, $\sqrt{\max\{(40/39)^3,40/38\}}$ bounds both
deletion factors. The binomial mode interval at $n-2$ is contained in
the wider interval at $38$. The positive inverse-index errors increase
when evaluated at $40$. A retained negative penalty is kept outside
this substitution: an expression $(zr+\mathrm{Dist})/n-zP$ uses
$zr+\mathrm{Dist}\ge0$ and keeps $-zP$ unchanged. For $d\ge0$,
$\tau>0$, its bad-pair factor satisfies
\[
 \frac{d}{n\tau+d}\le\frac{d}{40\tau+d}
 \le\frac{40\tau}{(40\tau+s)^2}d+
       \frac{s^2}{(40\tau+s)^2}\qquad(s\ge0),
\]
the tangent bound for a concave function. Thus the same actual
moment/noise triangle and every positive-width and signed derivative
gate license the confinement and its feedback at all three indices.

Here the numerical constants in the confinement proof of
Section~\ref{conf:re:confinement} are parameters. More precisely,
replace its floor $47$ by $N\ge40$, its initial strip $.55$ by a
number $w\in(0,1)$, and its conditional cubic coefficient by
$c_\rho>0$. Retain a probability panel $[l,h]\subset(0,1/2)$ and assume
\[
 \operatorname{supp}Y\subset(-w,1+w),\quad d_0<w,
 \quad c_\rho^2[l(1-l)]^3>[3(1-l)+w]^2.
\]
The strip and the inward bounds $U_0\le d_0$, $U_1\ge-d_0$
together give $|U_i|<w$. The displayed conditional cubic identity gives
$\rho\le\rho_B+\ell^3(3q+w)s\le\rho_B+c_\rho s$.
Use $J_0=N-2$ for the binomial mode interval and the deletion factor
$\sqrt{\max\{(N/(N-1))^3,N/(N-2)\}}$. In the interval geometry
use $g=1-w-2d_0>0$ and $r\le w+d_0$; in the upper contradiction
use $L_g-wU_+>\mathcal E$. The other width, disjointness,
positive-slope and budget gates are unchanged. These substitutions
follow directly through the proof: $w$ enters only the conditional
cubic bound and the endpoints of the known strip, while the index
enters only the decreasing factors already bounded above. They
therefore prove the same confinement implication whenever all the
resulting gates pass.

The central sequence uses $N=40$, $w=.55$ initially and
$c_\rho=24$. The left rows start with their separately acquired
strip $w\in\{.55,.575,.6\}$ and use $c_\rho=30$. After a strict
contraction, the proved amplitude becomes the next strip width.
Thus every later use has an already established strip, and the
left sequence does not assume the narrower uniform-tail strip.

\subsection{The full acquired domain and its local cover}

The initializer covers the whole moment range from $\rho=1$ to the
Shevtsova upper cap, not merely a low-moment subinterval. Its fourteen
refinements and final export each cover all signed second and third
moments and the common first moment; its probability table has 100
closed cells. The left acquisition covers its 38 cells below $.38$,
refines all twelve surviving cells into sixty children and excludes
the first eleven children. The remaining 49 rows cover $[.282,.38]$.
The central acquisition covers the other 62 coarse cells, excludes
every cell above $.48$, and refines $[.38,.48]$ into fifty children;
the final two children are excluded. Its 48 remaining rows cover
$[.38,.476]$. Thus the probability and moment complements are part of
this acquisition itself.

Each local row uses the same-law joint resource and both signs of the
raw threshold. Zero noise is covered by the Bernoulli theorem. At
positive noise, nearest-count and whole-count witnesses come from
the common exact identity~\eqref{ext:local-count-identity}.
The central phase cover is
\[
\begin{array}{c|c|c}
 \lambda&\text{nearest-count phase}&\text{whole-count phase}\\
 (0,.05]&\text{all}&--\\
 {[.05,.08]}&\text{all}&--\\
 {[.08,.10]}&|u|\ge.19&|u|\le.19\\
 {[.10,.12]}&|u|\ge.27&|u|\le.27\\
 {[.12,\Lambda]}&--&\text{all}.
\end{array}
\]
The left rows require a different cover:
\[
\begin{array}{c|c|c}
 \lambda&\text{nearest-count phase}&\text{whole-count phase}\\
 (0,.05]&\text{all}&--\\
 {[.05,.0925]}&\text{all}&--\\
 {[.0925,.10]}&|u|\ge.16&|u|\le.16\\
 {[.10,.12]}&|u|\ge.25&|u|\le.25\\
 {[.12,.135]}&u\le-.27&-.27\le u\le.5\\
 {[.135,\Lambda]}&--&\text{all}.
\end{array}
\]
Every band is intersected with the row's acquired $\Lambda$ and all
seams are closed. The actual computations further split noise into
intervals of width $1/400$ and phase into intervals of width $1/1000$
where needed. The full-count families have 70 central signed rows and
82 left signed rows; the other rows end below the respective starting
noise levels and are already covered by the nearest-count tests.

At these three finite indices the proof uses a positive lower bound
$\epsilon_p$ for $\min_{n=40,41,42}(\CE-C^{\rm bin}_{n,p})$,
uniform on its probability cell.
The complete binomial check covers $p\in[.28,.48]$, both discrepancy
signs and every count at each actual index. A signed inverse correction
$I$ is bounded by
\[
 I/n\le\max\{I/40,I/42\}.
\]
For normalized mass $b=\sqrt n\pi_k\in[b_0,1]$, a bracket $J$ of either sign is
bounded above by $\max\{b_0J,J\}$. These replacements explain why
this finite comparison is not an unrestricted floor-40 theorem.
They also give the finite-slack version of the whole-count test.
Suppose $S$ bounds $\rho_B(C-C_B)/\lambda$ on a closed noise
band $[a,b]$ with $a>0$, and $\rho_B\ge\rho_{B,l}>0$.
It is enough to check
\[
                  S<\rho_{B,l}\epsilon_p/b.
\]
If $S<0$, then $C<C_B$; otherwise multiplication by
$\lambda\le b$ gives $C<C_B+\epsilon_p\le\CE$.
The saved full-count formulas have a bracket affine in $\lambda$
plus a positive multiple of $(\alpha\lambda+\beta)^{-1}$,
and an inverse correction that is constant plus a nonnegative
multiple of $1/\lambda$. The two sign-sensitive maxima above
preserve convexity. Thus their values at the two noise endpoints
bound the entire band. The finite check may use positive binomial
slack even where strict comparison with $C_B$ itself is unavailable.
For the nearest-count comparisons, the entropy segment
$p\pm\delta/40$ contains the actual segment $p\pm\delta/n$.
All log-factorial, variance and ratio errors have positive numerators.
The signed Gaussian derivative factors have the forms
$1-A_n/n$ and $(1-\eta_n/n)^{-1}$, with
$0\le A_n\le A_{40}$ and $0\le\eta_n\le\eta_{40}<40$.
Thus the lower bound for the first and upper bound for the second
are valid at the floor. Their phase and threshold domains and the
gate $40pq>\delta^2$ are retained. A signed cubic correction is
minimized over its whole positive multiplier interval before it is
used in a lower repair bound; its sign is not discarded by changing
the index.

The finite binomial slack itself has a direct one-dimensional proof.
On a probability cell $[a,b]\subset(0,1/2)$ put
$F_j(p)=\Pr(\operatorname{Bin}(n,p)\le j)$ and
$z_j(p)=(j-np)/\sqrt{np(1-p)}$. Both functions decrease with $p$,
while $\sqrt{np(1-p)}/(1-2p(1-p))$ increases. Therefore both jump
sides of the normalized discrepancy are bounded by
\[
 \frac{\sqrt{nb(1-b)}}{1-2b(1-b)}
 \max\{F_j(a)-\Phi(z_j(b)),\,
        \Phi(z_j(a))-F_{j-1}(b),\,0\}.
\]
Here $F_{-1}=0$. Taking every $j=0,\ldots,n$ covers all thresholds,
because between jumps the distribution function is constant and
$\Phi$ increases. Each of the 100 probability panels is divided
into 128 closed subcells. The binomial probabilities at the endpoints
are exact rational numbers. On $|z|\le5$, an integrated degree-158
Taylor polynomial for $e^{-z^2/2}$ with its explicit next-term
integral remainder $|z|^{161}/(2^{80}80!161)$ encloses $\Phi$
after multiplication by $\phi(0)$; outside this interval the proved Mills
bound is $3\cdot10^{-7}$. Directed integer square roots and a rational
Machin enclosure for $\pi$ complete the enclosure. Subtracting these
upper scores from a rational lower bound for $\CE$ gives the saved
positive $\epsilon_p$ with both atom conventions included.

\paragraph{Binding the actual inputs.}
The comparison's amplitude profiles are the exact rational-string
acquisition object, renamed \texttt{profiles.json}; their complete bytes
must agree with \texttt{profile\_candidate.json}. The same source binding
is checked for the probability targets and for every actual-index
binomial table. Thus the numerical comparisons consume the data already
proved for the same marginal.

\paragraph{Numerical completion.}
The finite conclusion requires the ten fresh acquisition actions,
the separate indexed support check, all nine companion checks, all
152 full-count rows with their separate arithmetic replays and two
complete binders. The numerical
status of these branches is recorded separately; source identities and
historical success fields do not supply any of these inequalities.
